\answer{ex:04:1}{(а)~$\approx1.1\cdot10^{11}$ біт/Дж. (б)~$\log_2\sqrt{10}\approx1.7$ біта проти $81$: у $\approx50$ разів.}
\answer{ex:04:2}{$r^*=(P_0/E_{\text{імп}})\ln\bigl(1/(r^*\Delta t)\bigr)$, $P_0/E_{\text{імп}}=1.16\,\text{с}^{-1}$: $r^*\approx6$ імпульсів за секунду, $7.4\cdot10^{10}$ біт/Дж.}
\answer{ex:04:3}{$\mathcal I=\frac{N}{1+(N-1)c}=9.9$ (границя $10$) проти $\mathcal I=\frac{N}{1-c}=1111$: у $\approx110$ разів; кореляція шкідлива, лише коли схожа на сигнал.}
\answer{ex:04:4}{$\theta\approx68.7$ і $19.9$ (в одиницях $w$), $m=19$ і $m=10$: швидкому отримувачеві потрібно майже вдвічі менше збіжних входів, хоча його середній вхід уп'ятеро менший.}
\answer{ex:04:5}{$U\tau_{\text{відн}}\ge0.725\,\text{с}$ і $\tau_{\text{відн}}(1-2U)\le0.05\,\text{с}$, що вимагає $U\ge0.483$; найменше $\tau_{\text{відн}}=0.725\,\text{с}$ при $U=1$ ($1.45\,\text{с}$ при $U=0.5$).}
\answer{ex:04:6}{(а)~$\theta\,e/(e-1)\approx1.58\,\theta$ за будь-яких $\tau$ і $\theta$. (б)~$14$ імпульсів.}
\answer{ex:04:7}{(а)~$1.62$ біта на слово: $0.77$ --- кількість імпульсів, $0.84$ --- їхнє розташування. (б)~$1.13$ і $1.59$ біта: за $30$ словами оцінка втрачає майже третину.}
